
This section defines routines that shall be invoked from within \ac{GPU} device
kernels. Calling these routines from the host results in \textbf{undefined
behavior}.

The \openshmem library shall have been initialized with
\CONST{SHMEM\_DEVICE\_KERNEL\_INIT} and the corresponding capability shall have
been returned in the \VAR{provided} argument of \FUNC{shmem\_init\_thread}.
Calling any routine in this section without the required initialization results
in \textbf{undefined behavior}.

All routines in this section use the \texttt{shmemg\_} prefix and operate on
symmetric data objects allocated from the \ac{GPU} symmetric heap. A symmetric
address and a target \ac{PE} identifier are used to address remote data objects.
Unless otherwise noted, any \ac{GPU} thread within a device kernel may call these
routines.

In the synopses that follow, the \CONST{SHMEMG\_DEVICE} annotation
(\cref{subsec:shmemg_device}) marks each routine as callable from a \ac{GPU}
device kernel. It expands to the execution-space qualifier required by the
compilation environment, or to no tokens where none is required. The code
examples in this specification use CUDA syntax, such as \texttt{\_\_global\_\_}
kernels and vendor-specific kernel-launch syntax, for concreteness; equivalent
constructs apply in other \ac{GPU} programming environments.

